\begin{dang}{Sự trao đổi nhiệt}
	\Anlythuyet{\Binhluan{
			\begin{itemize}
				\item Xác định nhiệt lượng vật 1 trao đổi: $Q_1=m_1.c_1.(t-t_1)$
				\item Xác định nhiệt lượng vật 2 trao đổi: $Q_2=m_2.c_2.(t-t_2)$\\
				trong đó: $t_1,\;t_2$ là nhiệt độ ban đầu của hai vật, t là nhiệt độ của hai vật khi đạt trạng thái cân bằng nhiệt.
				\item ....
				\item Xác định nhiệt lượng cấp vào hệ Q
				\item Áp dụng điều kiện cân bằng nhiệt: 
				\begin{align*}
					Q_1+Q_2+...+Q_n=Q
				\end{align*}
			\end{itemize}
		}
		{Khi sử dụng cách làm này ta giữ nguyên dấu của Q.\\
			$Q>0$: Nhiệt lượng thu vào.\\
			$Q<0$: Nhiệt lượng toả ra}
		\Note{
			Nếu chỉ có các vật trong hệ trao đổi nhiệt với nhau thì $Q=0$.
	}}
	
\end{dang}